\section{Introduction}
\label{sec:intro}

A QR code stores one bit per module~\cite{iso18004}. Printing three binary layers in the
color channels of one symbol triples the data in the same area: each module shows one of
eight colors, and its red, green and blue bits belong to the three
layers~\cite{andre2014colour,blasinski2013pccc,yang2018hiq}. Polychrome codes are also
standardized, as JAB Code~\cite{iso23634}. A multiplexed QR symbol keeps the layout of a
standard symbol, so once the module colors are known each layer is read with standard
Reed-Solomon decoding.

Recovering those colors from a phone photograph is the difficult step. Illumination and
white balance shift every sample, overlapping colorant and sensor spectra mix the channels
within a module, and lens blur, ink spread and JPEG chroma subsampling mix each module with
its neighbors~\cite{yang2018hiq}. The last effect grows with density. Blur stays roughly
fixed in pixels while modules shrink, so in a dense symbol a larger share of each sample
comes from neighboring modules. Small module error rates also matter: a symbol fails as soon
as one Reed-Solomon block in one layer holds more errors than it can correct. In our
experiments two classifiers that differ by three points in module accuracy differ by more
than a factor of two in symbols decoded (Section~\ref{sec:results}).

Earlier decoders handle these effects in several ways. Blasinski et al.\ estimate
cross-channel mixing per image from pilot blocks or by expectation
maximization~\cite{blasinski2013pccc,bulan2012em}. HiQ trains color classifiers offline on
labelled photographs and learns cross-module interference terms with
them~\cite{yang2018hiq}. Bagherinia and Manduchi decode blurred patches by iterative joint
estimation~\cite{bagherinia2014blur}, and screen-to-camera systems train networks end to
end~\cite{tamang2021deepd2c,tran2021deeplight}. A trained decoder depends on the printers,
displays and phones in its training set. We look for a method that needs no training data,
runs in time linear in the number of modules, and holds up on dense symbols.

Our method fits each photograph on its own. The finder, alignment and timing patterns of a
multiplexed QR symbol have fixed colors, so they label a few hundred samples per photograph.
Before fitting, a closed-form inverse of a symmetric leak between edge-adjacent modules,
with one fixed fraction $\beta=0.06$, reduces cross-module interference. The corrected
samples are then classified by regularized discriminant analysis with shrink
$\lambda=0.5$, fitted to the labelled samples. On the public CUHK-CQRC dataset the method
decodes \HiqCodeChosen\% of localized symbols, against \HiqCodeRda\% without leak inversion
and \HiqCodeEuclid\% for nearest-centroid classification; on a new dataset of qrism symbols
it decodes \QrCodeChosen\%. On HiQ's own test split, our decoder reads \PairCommonOurs\% of
the symbols that both it and the released HiQ decoder localize, against \PairCommonHiq\% for
HiQ.

The contributions are:
\begin{itemize}
\item a training-free color recovery method for multiplexed QR codes: per-photograph
calibration from function patterns, a one-parameter inverse of inter-module leakage, and
regularized discriminant analysis (Section~\ref{sec:method});
\item a dataset of \QrAttempted{} photographs of qrism symbols at seven versions, four
error-correction levels and four phones, with the source render of every symbol;
\item an evaluation on two datasets in which both parameters are chosen on each dataset and
checked on the other, with paired tests, and an end-to-end comparison with the released HiQ
decoder on its test split that also corrects a defect in its CMI step
(Sections~\ref{sec:setup} and~\ref{sec:results});
\item an open-source Rust implementation. \todo{repository URL}
\end{itemize}

Section~\ref{sec:related} reviews related work, Section~\ref{sec:symbol} describes the
symbol designs and decoding pipeline, Section~\ref{sec:method} the method,
Sections~\ref{sec:setup} and~\ref{sec:results} the experiments, and
Section~\ref{sec:limitations} the limitations.
